\documentclass[10pt]{article}

\usepackage{amsmath,amssymb,amsthm}
\usepackage[a4paper,landscape,margin=0.35in,includehead]{geometry}
\usepackage{array}
\usepackage{booktabs}
\usepackage{enumitem}
\usepackage{microtype}
\usepackage{multicol}
\usepackage{titlesec}
\usepackage{fancyhdr}
\usepackage{draftwatermark}
\usepackage{atbegshi}
\usepackage[hidelinks]{hyperref}

% =========================
% 2-page hard cap
% =========================
\AtBeginShipout{%
  \ifnum\value{page}>2
    \AtBeginShipoutDiscard
  \fi
}

% =========================
% Watermark (optional)
% =========================
\SetWatermarkText{\href{https://qrsntu.org}{QRS@NTU \,|\, qrsntu.org}}
\SetWatermarkScale{0.9}
\SetWatermarkLightness{0.995}

% =========================
% Header (remove header rule line)
% =========================
\setlength{\headheight}{1pt}
\pagestyle{fancy}
\fancyhf{}
\fancyhead[L]{\normalsize \textbf{MH1201 Linear Algebra II - Cheat Sheet (Exam)}}
\fancyhead[R]{\normalsize\textbf{\href{https://qrsntu.org}{QRS@NTU \,|\, qrsntu.org}}}
\renewcommand{\headrulewidth}{0pt}
\renewcommand{\footrulewidth}{0pt}
\fancyfoot[C]{\tiny\thepage}

% =========================
% Tight typography
% =========================
\setlength{\parindent}{0pt}
\setlength{\parskip}{0pt}
\setlist{nosep,leftmargin=1.15em}
\setlength{\abovedisplayskip}{2pt}
\setlength{\belowdisplayskip}{2pt}
\setlength{\abovedisplayshortskip}{2pt}
\setlength{\belowdisplayshortskip}{2pt}
\setlength{\tabcolsep}{2.2pt}
\renewcommand{\arraystretch}{0.95}

% =========================
% Headings readability
% =========================
\titleformat{\section}
  {\bfseries\normalsize}
  {}
  {0pt}
  {\underline}
\titleformat{\subsection}
  {\bfseries\small}
  {}
  {0pt}
  {}
\titlespacing*{\section}{0pt}{3pt}{1pt}
\titlespacing*{\subsection}{0pt}{2pt}{0.5pt}

% =========================
% Quick macros
% =========================
\newcommand{\R}{\mathbb{R}}
\newcommand{\C}{\mathbb{C}}
\newcommand{\F}{\mathbb{F}}
\newcommand{\N}{\mathbb{N}}
\newcommand{\Z}{\mathbb{Z}}
\newcommand{\rank}{\operatorname{rank}}
\newcommand{\nullity}{\operatorname{nullity}}
\newcommand{\Ker}{\operatorname{ker}}
\newcommand{\Range}{\operatorname{range}}
\newcommand{\im}{\operatorname{im}}
\newcommand{\spann}{\operatorname{span}}
\newcommand{\tr}{\operatorname{tr}}
\newcommand{\diag}{\operatorname{diag}}
\newcommand{\proj}{\operatorname{proj}}

\begin{document}

% =========================
% PAGE 1
% =========================
\begin{multicols}{3}
\footnotesize
\raggedcolumns

\section*{Vector Spaces (Ch1)}

\subsection*{Definition}
A vector space $V$ over field $\F$ has operations $+:V\times V\to V$ and $\cdot:\F\times V\to V$ satisfying:\\
\textbf{Addition:} (A1) $u+v=v+u$; (A2) $(u+v)+w=u+(v+w)$; (A3) $\exists\bm{0}\in V$: $v+\bm{0}=v$; (A4) $\forall v\,\exists(-v)$: $v+(-v)=\bm{0}$.\\
\textbf{Scalar mult:} (M1) $1\cdot v=v$; (M2) $(ab)v=a(bv)$.\\
\textbf{Distributive:} (D1) $a(u+v)=au+av$; (D2) $(a+b)v=av+bv$.

\subsection*{Examples}
$\R^n$, $\C^n$, $\F^{m\times n}$ (matrices), $\F[x]$ (polynomials), $\F^\infty$ (sequences), $C(\R)$ (continuous functions), $\R^\R$ (all functions).

\subsection*{Subspaces}
$W\subseteq V$ is a subspace if:\\
(S1) $\bm{0}\in W$;\\
(S2) $u,v\in W\Rightarrow u+v\in W$;\\
(S3) $v\in W,\,\lambda\in\F\Rightarrow\lambda v\in W$.\\
\textbf{Kernel trick:} If $L:V\to W$ linear, then $\Ker(L)=\{v\in V:L(v)=\bm{0}\}$ is subspace of $V$.\\
\textbf{Union:} $W_1\cup W_2$ is subspace $\Leftrightarrow$ $W_1\subseteq W_2$ or $W_2\subseteq W_1$.\\
\textbf{Intersection:} $\bigcap_i W_i$ always subspace.\\
\textbf{Sum:} $W_1+W_2=\{w_1+w_2:w_1\in W_1,w_2\in W_2\}$ is subspace.

\subsection*{Direct sum}
$V=W_1\oplus W_2$ if $V=W_1+W_2$ and $W_1\cap W_2=\{\bm{0}\}$.\\
Equiv: every $v\in V$ unique decomposition $v=w_1+w_2$.

\section*{Finite-Dimensional Spaces (Ch2)}

\subsection*{Linear combinations}
$v=\sum_{i=1}^n a_iv_i$ is lin.\ comb.\ of $v_1,\ldots,v_n$.\\
$\spann\{v_1,\ldots,v_n\}=\{\text{all lin.\ comb.\ of }v_i\}$ is subspace.

\subsection*{Linear independence}
$v_1,\ldots,v_n$ lin.\ indep.\ if $\sum a_iv_i=\bm{0}\Rightarrow$ all $a_i=0$.\\
Equiv: no $v_i$ is lin.\ comb.\ of others.\\
$v_1,\ldots,v_n$ lin.\ dep.\ if $\exists$ nontrivial $\sum a_iv_i=\bm{0}$.

\subsection*{Basis \& dimension}
Basis of $V$: linearly independent spanning set.\\
Every basis of $V$ has same size = $\dim V$.\\
If $\dim V=n$, then:\\
\quad Any $n$ lin.\ indep.\ vectors form basis.\\
\quad Any $n$ vectors spanning $V$ form basis.\\
\quad Any lin.\ indep.\ set can be extended to basis.\\
\quad Any spanning set contains a basis.

\subsection*{Dimension formulas}
If $W_1,W_2$ subspaces of $V$:\\
$\dim(W_1+W_2)=\dim W_1+\dim W_2-\dim(W_1\cap W_2)$.\\
Direct sum: $\dim(W_1\oplus W_2)=\dim W_1+\dim W_2$.

\subsection*{Coordinates}
If $\beta=\{v_1,\ldots,v_n\}$ basis, $v=\sum a_iv_i$, then\\
$[v]_\beta=(a_1,\ldots,a_n)^T$ is coordinate vector w.r.t.\ $\beta$.

\section*{Linear Transformations (Ch3)}

\subsection*{Definition}
$T:V\to W$ is linear if:\\
(L1) $T(u+v)=T(u)+T(v)$;\\
(L2) $T(\lambda v)=\lambda T(v)$.\\
Equiv: $T(au+bv)=aT(u)+bT(v)$.\\
Always: $T(\bm{0})=\bm{0}$; $T(-v)=-T(v)$.

\subsection*{Kernel \& range}
$\Ker(T)=\{v\in V:T(v)=\bm{0}\}$ (null space).\\
$\Range(T)=\{T(v):v\in V\}=\im(T)$ (image).\\
Both are subspaces.\\
$\nullity(T)=\dim\Ker(T)$; $\rank(T)=\dim\Range(T)$.

\subsection*{Rank-Nullity Theorem}
If $T:V\to W$ linear, $\dim V<\infty$, then\\
$\dim V=\nullity(T)+\rank(T)$.

\subsection*{Injectivity \& surjectivity}
$T$ injective (1-1) $\Leftrightarrow$ $\Ker(T)=\{\bm{0}\}$.\\
$T$ surjective (onto) $\Leftrightarrow$ $\Range(T)=W$.\\
If $\dim V=\dim W<\infty$, then\\
\quad $T$ injective $\Leftrightarrow$ $T$ surjective $\Leftrightarrow$ $T$ bijective.

\subsection*{Isomorphism}
$T:V\to W$ isomorphism if $T$ linear bijection.\\
$V\cong W$ if $\exists$ isomorphism $V\to W$.\\
$V\cong W$ $\Leftrightarrow$ $\dim V=\dim W$ (finite-dim).\\
$V$ finite-dim, $\dim V=n$ $\Rightarrow$ $V\cong\F^n$.

\subsection*{Composition}
$(S\circ T)(v)=S(T(v))$.\\
If $S,T$ linear, then $S\circ T$ linear.\\
$\Ker(T)\subseteq\Ker(S\circ T)$; $\Range(S\circ T)\subseteq\Range(S)$.

\subsection*{Invertibility}
$T:V\to V$ invertible if $\exists T^{-1}:V\to V$ s.t.\ $T\circ T^{-1}=T^{-1}\circ T=I$.\\
$T$ invertible $\Leftrightarrow$ $T$ bijective $\Leftrightarrow$ $\Ker(T)=\{\bm{0}\}$ (finite-dim).

\section*{Matrix Representations (Ch4)}

\subsection*{Matrix of linear map}
$T:V\to W$ linear, $\beta=\{v_1,\ldots,v_n\}$ basis of $V$, $\gamma=\{w_1,\ldots,w_m\}$ basis of $W$.\\
$[T]_\beta^\gamma$ is $m\times n$ matrix with $j$-th column $=[T(v_j)]_\gamma$.\\
$[T(v)]_\gamma=[T]_\beta^\gamma[v]_\beta$.

\subsection*{Matrix operations}
$(S+T)$ linear: $[S+T]=[S]+[T]$.\\
$(cT)$ linear: $[cT]=c[T]$.\\
Composition: $[S\circ T]=[S][T]$.

\subsection*{Change of basis}
$\beta,\beta'$ bases of $V$. Change-of-basis matrix:\\
$P_{\beta\to\beta'}=\big[[v_1]_{\beta'}\cdots[v_n]_{\beta'}\big]$ where $\beta=\{v_1,\ldots,v_n\}$.\\
Then $[v]_{\beta'}=P_{\beta\to\beta'}[v]_\beta$ and $P_{\beta'\to\beta}=P_{\beta\to\beta'}^{-1}$.\\
If $T:V\to V$, bases $\beta,\beta'$ of $V$, then\\
$[T]_{\beta'}=P_{\beta\to\beta'}[T]_\beta P_{\beta'\to\beta}=P^{-1}[T]_\beta P$ where $P=P_{\beta\to\beta'}$.

\subsection*{Similar matrices}
$A,B$ similar ($A\sim B$) if $\exists$ invertible $P$: $B=P^{-1}AP$.\\
Similar matrices represent same linear map w.r.t.\ different bases.\\
$A\sim B\Rightarrow$ $\det A=\det B$, $\tr A=\tr B$, $\rank A=\rank B$.

\section*{Eigenvalues \& Eigenvectors (Ch5)}

\subsection*{Definition}
$T:V\to V$ linear. $\lambda\in\F$ is eigenvalue if $\exists v\neq\bm{0}$: $T(v)=\lambda v$.\\
Such $v$ is eigenvector for $\lambda$.\\
Eigenspace: $E_\lambda=\{v\in V:T(v)=\lambda v\}=\Ker(T-\lambda I)$.

\subsection*{Characteristic polynomial}
$p_A(t)=\det(tI-A)$ is char.\ poly.\ of $A$.\\
$\lambda$ eigenvalue $\Leftrightarrow$ $p_A(\lambda)=0$.\\
Algebraic multiplicity $m_\lambda$: multiplicity of $\lambda$ as root of $p_A$.\\
Geometric multiplicity $g_\lambda=\dim E_\lambda$.\\
Always: $1\le g_\lambda\le m_\lambda$.

\subsection*{Eigenvalue properties}
$\det A=\prod\lambda_i$ (product of eigenvalues).\\
$\tr A=\sum\lambda_i$ (sum of eigenvalues).\\
Eigenvectors for distinct eigenvalues are lin.\ indep.

\subsection*{Diagonalization}
$A$ diagonalizable if $A\sim D$ (diagonal).\\
Equiv: $\exists$ invertible $P$: $P^{-1}AP=D=\diag(\lambda_1,\ldots,\lambda_n)$.\\
Columns of $P$ are eigenvectors; diagonal entries of $D$ are eigenvalues.\\
$A$ diagonalizable $\Leftrightarrow$ $V$ has basis of eigenvectors $\Leftrightarrow$ $\sum g_{\lambda_i}=n$.\\
Sufficient: all eigenvalues distinct.\\
If $A$ diagonalizable, then $A^k=PD^kP^{-1}$ where $D^k=\diag(\lambda_1^k,\ldots,\lambda_n^k)$.

\subsection*{Minimal polynomial}
$m_A(t)$ is minimal poly.\ of $A$: monic poly.\ of least degree s.t.\ $m_A(A)=0$.\\
$m_A(t)$ divides $p_A(t)$.\\
$A$ diagonalizable $\Leftrightarrow$ $m_A(t)$ splits with no repeated roots.

\subsection*{Jordan Canonical Form}
Every matrix similar to Jordan form $J=\diag(J_1,\ldots,J_k)$ where\\
$J_i=\begin{pmatrix}\lambda_i&1&0&\cdots&0\\0&\lambda_i&1&\cdots&0\\\vdots&\vdots&\ddots&\ddots&\vdots\\0&0&0&\lambda_i&1\\0&0&0&0&\lambda_i\end{pmatrix}$ (Jordan block).\\
JCF unique up to block order.\\
$A$ diagonalizable $\Leftrightarrow$ all Jordan blocks are $1\times1$.

\columnbreak

\section*{Inner Product Spaces (Ch6)}

\subsection*{Inner product}
$\langle\cdot,\cdot\rangle:V\times V\to\F$ is inner product if:\\
(IP1) $\langle v,v\rangle\ge0$; $\langle v,v\rangle=0\Leftrightarrow v=\bm{0}$ (pos.\ def.).\\
(IP2) $\langle u,v\rangle=\overline{\langle v,u\rangle}$ (conjugate symmetry).\\
(IP3) $\langle au+bv,w\rangle=a\langle u,w\rangle+b\langle v,w\rangle$ (linearity in 1st arg).\\
Standard: $\langle x,y\rangle=\sum x_i\overline{y_i}$ on $\F^n$.

\subsection*{Norm}
$\|v\|=\sqrt{\langle v,v\rangle}$.\\
$\|v\|\ge0$; $\|v\|=0\Leftrightarrow v=\bm{0}$.\\
$\|cv\|=|c|\|v\|$.\\
Triangle inequality: $\|u+v\|\le\|u\|+\|v\|$.\\
Distance: $d(u,v)=\|u-v\|$.

\subsection*{Cauchy-Schwarz Inequality}
$|\langle u,v\rangle|\le\|u\|\|v\|$.\\
Equality $\Leftrightarrow$ $u,v$ linearly dependent.

\subsection*{Orthogonality}
$u\perp v$ if $\langle u,v\rangle=0$.\\
$u\perp W$ if $u\perp w$ for all $w\in W$.\\
Orthogonal complement: $W^\perp=\{v\in V:v\perp W\}$.\\
$(W^\perp)^\perp=W$ (finite-dim); $\dim W+\dim W^\perp=\dim V$.\\
$V=W\oplus W^\perp$ (finite-dim).

\subsection*{Orthogonal sets}
$\{v_1,\ldots,v_n\}$ orthogonal if $\langle v_i,v_j\rangle=0$ for $i\neq j$.\\
Orthonormal if orthogonal and $\|v_i\|=1$ for all $i$.\\
Orthogonal sets are linearly independent.\\
Orthonormal basis: basis that is orthonormal.

\subsection*{Orthogonal projection}
$v\in V$, $W$ subspace with orthonormal basis $\{u_1,\ldots,u_k\}$.\\
$\proj_W(v)=\sum_{i=1}^k\langle v,u_i\rangle u_i$.\\
$v=\proj_W(v)+\proj_{W^\perp}(v)$ (orthogonal decomposition).\\
$\proj_W(v)$ is closest vector in $W$ to $v$: $\|v-\proj_W(v)\|=\min_{w\in W}\|v-w\|$.

\subsection*{Gram-Schmidt Orthogonalization}
Input: lin.\ indep.\ $\{v_1,\ldots,v_n\}$. Output: orthogonal $\{u_1,\ldots,u_n\}$ with $\spann\{u_1,\ldots,u_k\}=\spann\{v_1,\ldots,v_k\}$.\\
$u_1=v_1$;\\
$u_2=v_2-\frac{\langle v_2,u_1\rangle}{\langle u_1,u_1\rangle}u_1$;\\
$u_k=v_k-\sum_{i=1}^{k-1}\frac{\langle v_k,u_i\rangle}{\langle u_i,u_i\rangle}u_i$.\\
Normalize: $e_i=\frac{u_i}{\|u_i\|}$ to get orthonormal basis.

\subsection*{Orthogonal matrices}
$Q$ orthogonal if $Q^TQ=I$ (equiv: columns orthonormal).\\
$Q$ orthogonal $\Rightarrow$ $\|Qx\|=\|x\|$ (preserves norm).\\
$Q$ orthogonal $\Rightarrow$ $\det Q=\pm1$.

\section*{Quick Reference}

\subsection*{Common spaces}
$\R^n$: $\dim=n$; standard basis $\{e_1,\ldots,e_n\}$.\\
$P_n(\R)$ (poly.\ deg $\le n$): $\dim=n+1$; basis $\{1,x,\ldots,x^n\}$.\\
$\F^{m\times n}$: $\dim=mn$; basis $\{E_{ij}\}$ ($E_{ij}$ has 1 at $(i,j)$, 0 elsewhere).

\subsection*{Matrix rank formulas}
$\rank(A)=\rank(A^T)$.\\
$\rank(AB)\le\min\{\rank(A),\rank(B)\}$.\\
$\nullity(A)=n-\rank(A)$ ($A$ is $m\times n$).

\subsection*{Invertibility (square $n\times n$ matrix $A$)}
TFAE: (1) $A$ invertible; (2) $\det A\neq0$; (3) $\rank A=n$; (4) $\nullity A=0$; (5) cols of $A$ lin.\ indep.; (6) cols of $A$ span $\F^n$; (7) rows of $A$ lin.\ indep.; (8) $Ax=\bm{0}\Rightarrow x=\bm{0}$; (9) $Ax=b$ has unique solution for all $b$.

\subsection*{Determinant properties}
$\det(AB)=\det(A)\det(B)$.\\
$\det(A^T)=\det(A)$.\\
$\det(A^{-1})=1/\det(A)$.\\
$\det(cA)=c^n\det(A)$ ($A$ is $n\times n$).

\subsection*{Trace properties}
$\tr(A+B)=\tr(A)+\tr(B)$.\\
$\tr(cA)=c\tr(A)$.\\
$\tr(AB)=\tr(BA)$ (cyclic property).\\
$\tr(A)=\sum\lambda_i$ (eigenvalues).

\subsection*{Block matrices}
$\det\begin{pmatrix}A&B\\0&D\end{pmatrix}=\det(A)\det(D)$ (block triangular).

\subsection*{Quick checks}
Not subspace: zero vector missing OR not closed under $+$ OR not closed under $\cdot$.\\
Lin.\ dep.: one vector is lin.\ comb.\ of others OR more vectors than dimension.\\
Not diagonalizable: $g_\lambda<m_\lambda$ for some $\lambda$ OR repeated eigenvalue with deficient eigenspace.

\subsection*{Exam tips}
\textbf{Subspace:} Check (S1), (S2), (S3) OR express as kernel OR parametrize and take span.\\
\textbf{Basis:} Show lin.\ indep.\ AND spanning OR show correct size and lin.\ indep.\ OR correct size and spanning.\\
\textbf{Lin.\ indep.:} Set $\sum a_iv_i=\bm{0}$, solve, conclude $a_i=0$.\\
\textbf{Dimension:} Count basis size OR use dimension formula.\\
\textbf{Kernel:} Solve $T(v)=\bm{0}$ OR $Av=\bm{0}$; express as span.\\
\textbf{Range:} Find $T(v_i)$ for basis vectors $v_i$; take span.\\
\textbf{Eigenvalues:} Solve $\det(tI-A)=0$.\\
\textbf{Eigenvectors:} Solve $(A-\lambda I)v=\bm{0}$ for each $\lambda$.\\
\textbf{Diagonalization:} Check if $\sum g_{\lambda_i}=n$; form $P$ from eigenvectors.\\
\textbf{Gram-Schmidt:} Apply formula sequentially; normalize at end if needed.\\
\textbf{Orthogonal proj.:} Use $\proj_W(v)=\sum\langle v,u_i\rangle u_i$ with orthonormal $\{u_i\}$.

\end{multicols}

\end{document}